\section{Introduction}

Hospitals record every diagnosis a patient receives in an electronic health record (EHR). One natural way to study these records is as a network: here each node is a hospital stay, and an edge connects two stays whose diagnosis histories look alike. The idea is the same as a film recommender that connects people with similar viewing histories. If similar viewers liked a film, you might too; if similar patients had a long hospital stay, this patient might too.

We ask two questions. First, \emph{are networks of similar patients organised around clinical outcomes, and what structure do they have?} Second, \emph{what changes if the network accidentally uses information from the future?} This second problem is called \emph{information leakage}. A network built from diagnoses recorded after the moment we are trying to predict can look informative only because it has already seen the answer. Leakage is a common and serious mistake in medical machine learning, because it makes methods look better on paper than they could be in a real hospital.

We answer both questions with classical network analysis and no machine learning, so that every result can be traced to a simple, transparent measure.
